\chapter{Model}

\section{Markowskie procesy decyzyjne}

Model markowskiego procesu decyzyjnego (MDP) z dyskontowaniem quasi-hiperbolicznym jest zdefiniowany jako: $M = (S, A, q, r, \alpha, \beta)$, gdzie:

\begin{itemize}
\item $S$ -- zbiór stanów,
\item $A$ -- zbiór akcji,
\item $q(s' \mid s,a)$ -- prawdopodobieństwo przejścia ze stanu $s$ do stanu $s'$ przy akcji $a$,
\item $r : S \times A \rightarrow \mathbb{R}$ -- funkcja nagrody,
\item $\alpha > 0$ -- parametr wzmacniający wagę nagród natychmiastowych,
\item $\beta \in [0,1)$ -- standardowy współczynnik dyskontowania wykładniczego.
\end{itemize}

\subsection{Przestrzenie pomocnicze}

Oznaczmy przez:
\begin{enumerate}
\item $\Pr(A)$ -- zbiór miar probabilistycznych na $A$; dla skończonego $A$ utożsamiamy z $\mathbb{R}^{|A|}$.
\item $\mathcal{H}_n$ -- historie długości $n$: $h_n = (s_0, a_0, s_1, a_1, \ldots, s_n)$, przy czym $h_0 = s_0$.
\end{enumerate}

Historia gromadzi pełen kontekst interakcji agent--środowisko do czasu $n$.

\subsection{Polityki decyzyjne}

\begin{enumerate}
\item Ogólna polityka: $\phi = (\phi_0, \phi_1, \ldots)$, gdzie $\phi_n : \mathcal{H}_n \rightarrow \Pr(A)$ określa rozkład akcji dla każdej historii $h_n \in \mathcal{H}_n$, $n \in \mathbb{N} \cup \{0\} = \mathbb{N}_0$.
\item Polityki markowskie: zależne od stanu i czasu, $\phi_n : S \rightarrow \Pr(A)$.
\item Polityki stacjonarne: zależne tylko od stanu, $\phi_s : S \rightarrow \Pr(A)$; będziemy oznaczać je przez $\phi_s$.
\end{enumerate}

Zbiory polityk:
\begin{itemize}
\item $\Phi$ -- zbiór wszystkich polityk zależnych od historii,
\item $\Phi_m$ -- zbiór polityk markowskich,
\item $\Phi_s$ -- zbiór polityk stacjonarnych.
\end{itemize}

\subsection{Funkcja wypłaty}

Rozważamy oczekiwaną wypłatę w nieskończonym horyzoncie czasowym z dyskontowaniem quasi-hiperbolicznym. Współczynniki w kolejnych okresach wynoszą 1, $\alpha \beta$, $\alpha \beta^2$, $\alpha \beta^3$, \ldots.

Wypłata dyskontowana w nieskończonym horyzoncie czasowym dla polityki $\phi \in \Phi$ i stanu początkowego $s_0$ jest dana wyrażeniem
\begin{align}
    J_{\phi}^{\alpha,\beta}(s_0)
    &= \mathbb{E}_{\substack{a_0 \sim \phi_0(\cdot \mid s_0) \\ s' \sim q(\cdot \mid s_0,a_0)}}
      \left[ r(s_0,a_0) + \alpha\,\beta\, J_\phi^{\beta}(s') \right] \nonumber \\
    &= \sum_a \phi_0(a \mid s_0) \left[ r(s_0,a) + \alpha\,\beta\, \sum_{s'} q(s' \mid s_0,a) J_\phi^{\beta}(s') \right],
\label{eq:qh-return}
\end{align}
gdzie $\overrightarrow{\phi} = (\phi_1, \phi_2, \ldots)$, a $J_{\phi}^{\beta}(s')$ oznacza standardową wypłatę oczekiwaną z dyskontowaniem wykładniczym
\begin{equation}
J_\phi^{\beta}(s') = \sum_{n=1}^{\infty} \beta^{n-1} \; r^{(n)}(s', \overrightarrow{\phi}).
\label{eq:beta-value}
\end{equation}

\section{Quasi-hiperboliczne dyskontowanie}

Dyskontowanie quasi-hiperboliczne jest zadane przez: $d(0) = 1$, a dla $t \ge 1$ mamy $d(t) = \alpha\,\beta^t$.

W przeciwieństwie do wykładniczego dyskontowania, gdzie $d(t) = \gamma^t$, dyskontowanie quasi-hiperboliczne lepiej modeluje preferencje niespójne w czasie. Parametr $\alpha$ wprowadza uprzedzenie teraźniejszości (ang. \textit{present bias}), co powoduje, że decydent relatywnie bardziej ceni nagrody natychmiastowe w porównaniu z przyszłymi.

\subsection{Przykład z jabłkiem}

Rozważmy prosty przykład decyzji dotyczącej konsumpcji jabłka. Załóżmy, że mamy następującą sytuację:

\begin{itemize}
\item W dniu $t=0$ możemy wybrać między zjedzeniem małego jabłka dziś lub dużego jabłka za tydzień.
\item Małe jabłko daje zadowolenie o wartości $r_1 = 5$ jednostek.
\item Duże jabłko daje zadowolenie o wartości $r_2 = 10$ jednostek.
\item Parametry dyskontowania: $\alpha = 0.8$, $\beta = 0.95$.
\end{itemize}

\textbf{Przy dyskontowaniu wykładniczym} ($d(t) = \beta^t$):
\begin{itemize}
\item Wartość małego jabłka dzisiaj: $5 \cdot \beta^0 = 5$
\item Wartość dużego jabłka za tydzień: $10 \cdot \beta^7 \approx 6.98$
\item Decydent wybiera duże jabłko za tydzień.
\end{itemize}

\textbf{Przy dyskontowaniu quasi-hiperbolicznym} ($d(0) = 1$, $d(t) = \alpha\beta^t$ dla $t \geq 1$):
\begin{itemize}
\item Wartość małego jabłka dzisiaj: $5 \cdot 1 = 5$
\item Wartość dużego jabłka za tydzień: $10 \cdot \alpha\beta^7 = 10 \cdot 0.8 \cdot 0.95^7 \approx 5.58$
\item Decydent może wybrać duże jabłko, ale wybór jest mniej oczywisty.
\end{itemize}

\textbf{Niespójność czasowa:} Załóżmy, że minął tydzień i teraz stoimy przed wyborem między małym jabłkiem dziś a dużym za kolejny tydzień. Przy dyskontowaniu quasi-hiperbolicznym ponownie obliczamy:
\begin{itemize}
\item Wartość małego jabłka dzisiaj: $5 \cdot 1 = 5$
\item Wartość dużego jabłka za tydzień: $10 \cdot 0.8 \cdot 0.95^7 \approx 5.58$
\end{itemize}

Zauważmy, że relatywne wartości są takie same jak tydzień temu, ponieważ uprzedzenie teraźniejszości $\alpha$ zawsze faworyzuje nagrodę natychmiastową.

\subsection{Przykład z podjęciem decyzji natychmiastowego zakupu}

Rozważmy decyzję zakupową w sklepie internetowym:

\begin{itemize}
\item Produkt kosztuje 1000 zł.
\item Mamy dwie opcje:
  \begin{itemize}
  \item Opcja A: Kupić teraz i zapłacić pełną cenę 1000 zł.
  \item Opcja B: Czekać miesiąc na promocję i zapłacić 800 zł.
  \end{itemize}
\item Użyteczność z posiadania produktu: 100 jednostek na miesiąc.
\item Koszt pieniężny: utrata użyteczności proporcjonalna do wydanej kwoty.
\end{itemize}

\textbf{Analiza decyzji:}

Niech $\alpha = 0.7$, $\beta = 0.98$ (dyskontowanie dzienne, przeskalowane na miesiąc: $\beta^{30} \approx 0.545$).

\textbf{Opcja A (kupno natychmiastowe):}
\begin{itemize}
\item Koszt: $-1000$ jednostek dzisiaj
\item Korzyść: $100$ jednostek miesięcznie przez wiele miesięcy
\item Wartość netto w pierwszym miesiącu: $-1000 + 100 = -900$
\item Wartość w następnych miesiącach: $\alpha \beta^{30} \cdot 100 + \alpha \beta^{60} \cdot 100 + \ldots$
\end{itemize}

\textbf{Opcja B (czekaj na promocję):}
\begin{itemize}
\item Koszt za miesiąc: $\alpha \beta^{30} \cdot (-800) \approx -304$
\item Korzyść od drugiego miesiąca: $\alpha \beta^{60} \cdot 100 + \alpha \beta^{90} \cdot 100 + \ldots$
\end{itemize}

Decydent z quasi-hiperbolicznym dyskontowaniem może wybrać opcję A (kupno natychmiastowe), mimo że opcja B jest ekonomicznie bardziej korzystna, ponieważ uprzedzenie teraźniejszości $\alpha < 1$ znacząco zmniejsza wartość przyszłych korzyści.

\textbf{Implikacje:}
\begin{itemize}
\item Decydent wykładniczy ($\alpha = 1$) wybrałby racjonalnie opcję B.
\item Decydent quasi-hiperboliczny ($\alpha < 1$) może ulec impulsowi i wybrać opcję A.
\item To wyjaśnia zjawiska takie jak zakupy impulsowe, nadmierne korzystanie z kart kredytowych, czy trudności z oszczędzaniem.
\end{itemize}
